%% chapter5.tex — Conclusions and Future Work

\chapter{Conclusions and Future Work}
\label{chap:conclusions}

This dissertation set out to add dynamic, multi-hop routing to the RA-TDMAs+
framework while preserving its TDMA slot discipline, and to do so in a way that
is reliable and observable enough for a real teleoperation link. The system
developed here meets that goal: it routes at Layer~3 over a TUN overlay, carries
the video on a TCP data plane separated from the UDP control beacons, and adds
per-link quality monitoring and a per-slot byte budget to the routing and
transmission path. The experimental evaluation of Chapter~\ref{chap:evaluation}
showed that the control channel is responsive and loss-free, that the video is
delivered at full rate, that recovery through the relay after a link failure is
automatic, and, centrally, that the nodes keep to their TDMA slots for both
periodic beacons and continuous video, with the relay leaving latency,
throughput and timing essentially unchanged.

Taken together, these results show that the system is a robust and fully
functional alternative to the Layer~2 predecessor of
Morais~\cite{morais2024synchronization}. It follows a different architectural
route, routing at Layer~3 with kernel forwarding over a TCP data plane, and it
was proven on real hardware to match that approach at the TDMA level while
gaining reliable in-order delivery, standard observability, and a relayed stream
that stayed stable across every experiment. It is, in short, a different and more
robust way of solving the same problem.

The work also exposed the limits of the chosen design, and these
define the most useful directions for future work. Two stand out. The first is
tighter control over \emph{when} data is actually sent. The second is a stronger
convergence guarantee for the case in which the lightweight beacons outlive a
broken TCP data plane. Both would be reinforced by a physical-layer link-quality
signal such as RSSI, which the current hardware did not allow us to use.


%% -----------------------------------------------------------------------
\section{Future Work}
\label{sec:future-work}

\subsection{Wire-Level Slot Containment}
\label{subsec:fw-slot-containment}

The framework controls how \emph{much} data it offers per slot, through the
per-slot byte budget and back-pressure mechanism (Section~\ref{sec:slot-control}),
but not the exact transmit instant. Because the data plane runs over TCP, a
\texttt{send()} only hands the data to the kernel, which then schedules
transmission on its own terms (TCP flow and congestion control, then CSMA/CA
medium access); a batch offered near the slot end can therefore spill slightly
into the next slot, the small, bounded incursion observed in
Section~\ref{subsubsec:video-slot}. The byte budget keeps this localised, but it
is not a hard, wire-level guarantee.

Several refinements would tighten it, from least to most invasive: a
conservative, empirically calibrated budget that leaves the kernel headroom to
drain before the boundary; intra-slot pacing of the packets across the slot
window; bounding the kernel's unsent backlog (\texttt{TCP\_NOTSENT\_LOWAT} and a
small \texttt{SO\_SNDBUF}); kernel-assisted pacing (an \texttt{fq} qdisc with
\texttt{SO\_MAX\_PACING\_RATE}); and, most thoroughly, gating transmission below
the socket by sending the slot-critical traffic as raw \texttt{AF\_PACKET} frames
timed to the slot. The last recovers the wire-level timing that the
connectionless Layer~2 design of Morais~\cite{morais2024synchronization} has by
construction, at the cost of re-introducing some of the complexity the TCP data
plane was chosen to avoid. Quantifying how much each refinement reduces the
incursions, with on-air measurement, is left as a direct continuation of this
work.

\subsection{Faster Convergence}
\label{subsec:fw-convergence}

Recovery after a link failure is dominated by the node-aging timeout
(\texttt{MAX\_AGE} $= 2.0$\,s) used to detect a dead direct link
(Section~\ref{subsec:res-convergence}); the routing computation itself is
negligible. Reducing \texttt{MAX\_AGE}, or replacing the fixed timeout with an
adaptive scheme that reacts to the recent beacon-loss rate, would speed up
recovery, traded against robustness to transient beacon loss in the shared
collision domain.

A more fundamental gap concerns a link that degrades unevenly across the two
planes: the lightweight UDP beacons can keep arriving after the heavier TCP data
plane has already broken (Section~\ref{subsec:res-demo}), so the topology still
believes the direct link is healthy while the operator has already lost the video
and the command channel, and the system does not reroute. Closing this gap
requires a \emph{data-plane liveness signal}: treating a stalled or failed TCP
connection to the next hop, like a missed beacon, as evidence to suppress the
link and force a reroute. This was attempted here, marking a link dead once its
TCP connection died and triggering an immediate reroute, but it did not yield
stable recovery on the testbed and is left as future work. The dominant cost to
cut is the source-side wind-down: after the failure the source can stall on the
dead socket up to the one-second send timeout (\texttt{SO\_SNDTIMEO}) before the
stream settles onto the relay, and a liveness signal is the most direct way to
shorten that tail.

\subsection{Physical-Layer Link Quality and RSSI}
\label{subsec:fw-rssi}

This work already routes on a link-quality metric, the per-link packet delivery
ratio (PDR) derived from beacon-reception statistics
(Section~\ref{subsec:res-slots}), which drives both the MST edge weights and
relay activation. PDR is a robust, protocol-level signal, but it is not the most
effective one, because it is lagging: it registers a failing link only after
enough beacons have already been lost. A physical-layer signal such as the
received signal strength indicator (RSSI) would complement it by reacting before
packets are lost, letting the MST shift weight toward a stronger path
pre-emptively rather than detecting the break after the fact. Incorporating RSSI
was attempted but proved impractical: the Raspberry~Pi adapters do not expose a
reliable per-neighbour RSSI in ad-hoc (IBSS) mode, and RSSI is in any case noisy
and would have to be filtered and combined with PDR rather than trusted alone.
With suitable hardware, an RSSI-and-PDR composite weight on the MST, together
with the data-plane liveness signal above, would give the system a faster and
more certain reaction to the link degradations that matter most for a mobile
teleoperation link.

\subsection{Multi-Hop Scalability}
\label{subsec:fw-multihop}

The evaluation used a three-node testbed, the minimum that exercises a relay. The
forwarding step is identical at every relay (a kernel \texttt{ip\_forward} of the
unmodified IP datagram, Section~\ref{sec:pkt-flow}), so the design generalises to
longer paths without change. What grows with the hop count is the per-hop slot
latency and the capacity sharing inherent to a single-channel multi-hop chain. A
systematic study on a larger testbed, ideally with nodes that cannot all hear one
another, would characterise these effects and separate them from the artefacts of
a single shared collision domain.

\subsection{Closing the Layer~2 Comparison}
\label{subsec:fw-l2}

The side-by-side re-implementation of the Layer~2 method
(Section~\ref{subsec:l2-rerun}) was reconstructed from the published description
rather than the original source. Separating how much of the streaming instability
it exhibited is intrinsic to the non-slotted relay, how much is an artefact of the
shared-domain testbed and the software-induced link failure, and how much is a
limitation of the reconstruction, would put the comparison on an even firmer
footing.
